\documentclass[10pt,a4paper]{amsart}
\usepackage{amsmath,amssymb,amsthm,graphicx}
\usepackage{amsmath,amssymb,amsthm,xfrac,xspace,hyperref,graphicx,comment}
\excludecomment{DefTralics}

\providecommand{\bibliofont}{\small}
\newcommand{\ie}{i.e.,\xspace}
\newcommand{\eg}{e.g.,\xspace}
\newcommand{\etc}{etc.\xspace}
\newcommand{\cf}{cf.\xspace}
\newcommand{\resp}{resp.\xspace}
\providecommand{\smaller}{\small}
\newcommand{\Subsection}{\subsection}
\newcommand{\Subsubsection}{\subsubsection}
\newcommand{\skpt}{\smallskip}
\emergencystretch=3em
\numberwithin{equation}{section}\providecommand{\No}{}
\usepackage{mathrsfs}
\let\mathcal\mathscr
\usepackage{mathtools}
\newcommand{\Changel}{\mathcode`l="7160}
\newcommand{\Changelback}{\mathcode`l="716C}
\newcommand{\NoChangel}[1]{%
\expandafter\let\csname old\string#1\endcsname=#1
\let#1=\relax
\newcommand{#1}{\mathcode`l="716C\csname old\string#1\endcsname\mathcode`l="7160 }%
}
\def\mto{\mathchoice{\longmapsto}{\mapsto}{\mapsto}{\mapsto}}
\def\hto{\mathrel{\lhook\joinrel\to}}
\def\ltwoheadrightarrow{\longrightarrow\hspace*{-5mm}\longrightarrow}
\def\stwoheadrightarrow{\rightarrow\hspace*{-3mm}\rightarrow}
\def\twoheadrightarrow{\mathchoice{\ltwoheadrightarrow}{\stwoheadrightarrow}{\stwoheadrightarrow}{\stwoheadrightarrow}}
\def\To#1{\mathchoice{\xrightarrow{\textstyle\kern4pt#1\kern3pt}}{\stackrel{#1}{\longrightarrow}}{}{}}
\def\Hto#1{\mathrel{\lhook\joinrel\To{#1}}}
\newcommand{\sbullet}{{\scriptscriptstyle\bullet}}

\NoChangel{\log}
\NoChangel{\ln}
\NoChangel{\lim}
\NoChangel{\limsup}
\NoChangel{\liminf}
\NoChangel{\varlimsup}
\NoChangel{\varliminf}
\NoChangel{\varinjlim}
\NoChangel{\varprojlim}

\let\dpl\displaystyle
\let\ts\textstyle

\AtBeginDocument{%
\def\labelitemi{$\sbullet$}%
}

\usepackage{graphicx}
\usepackage{multirow}

\usepackage[noabbrev,capitalize]{cleveref}
\crefdefaultlabelformat{#2\textup{#1}#3}

\usepackage{tikz}
\usepackage{tikz-cd}
\tikzset{math3d/.style=
{x= {(-0.353cm,-0.353cm)}, z={(0cm,1cm)},y={(1cm,0cm)}}}
\tikzset{JLL3d/.style=
{x= {(0.4cm,-0.2cm)}, z={(0cm,1cm)},y={(-1cm,0cm)}}}

\theoremstyle{plain}
\newtheorem{proposition}{Proposition}
\newtheorem{theorem}{Theorem}
\newtheorem*{problem}{Problem}
\newtheorem{corollary}{Corollary}
\newtheorem{lemma}{Lemma}
\newtheorem{conjecture}{Conjecture}

\theoremstyle{definition}
\newtheorem{definition}{Definition}
\newtheorem{remark}{Remark}
\newtheorem{example}{Example}

\newcommand{\RR}{\mathbb{R}}
\newcommand{\ZZ}{\mathbb{Z}}
\newcommand{\ZP}{\mathbb{Z}_{>0}}
\newcommand{\K}{{K}}
\newcommand{\La}{\mathcal{L}}
\newcommand{\Tam}[1]{\mathrm{PBT}_{#1}}
\newcommand{\PT}[1]{\mathrm{PT}_{#1}}
\newcommand{\sC}{\mathcal{C}}
\newcommand{\PolySub}{\mathsf{Poly}}\NoChangel{\PolySub}
\newcommand{\Top}{\mathsf{Top}}
\newcommand{\CW}{\mathsf{CW}}
\newcommand{\sF}{\mathcal{F}}
\DeclareMathOperator{\tp}{top}
\DeclareMathOperator{\bm}{bot}
\DeclareMathOperator{\pr}{pr}
\DeclareMathOperator{\conv}{conv}
\DeclareMathOperator{\diam}{diam}
\newcommand{\ang}[1]{\langle #1\rangle}
\renewcommand{\No}{\mathcal{N}}
\newcommand{\Ve}{\mathcal{V}}
\newcommand{\Bool}[1]{\{0, 1\}^{#1}}
\newcommand{\sk}{\mathrm{sk}}
\DeclareMathOperator{\vol}{vol}\NoChangel{\vol}
\newcommand{\tr}{\mathrm{tr}}
\newcommand{\id}{\mathrm{id}}
\newcommand{\itd}[1]{\triangle^{(#1)}}
\DeclareMathOperator*{\colim}{\mathrm{colim}}
\DeclareMathOperator{\im}{\mathrm{Im}}
\newcommand{\abs}[1]{\lvert #1\rvert}
\DeclareMathOperator{\codim}{codim}

\newcommand{\TreeR}
{\vcenter{\hbox{
\begin{tikzpicture}[yscale=0.2,xscale=0.2]
\draw[thick] (0,-1)--(0,0) -- (2,2);
\draw[thick] (1,1)--(0,2) ;
\draw[thick] (0,0)--(-1,1) ;
\draw [fill] (0,-1) circle [radius=0.035];
\draw [fill] (2,2) circle [radius=0.035];
\draw [fill] (0,2) circle [radius=0.035];
\draw [fill] (-1,1) circle [radius=0.035];
\end{tikzpicture}}}}
\newcommand{\TreeL}
{\vcenter{\hbox{
\begin{tikzpicture}[yscale=0.2,xscale=0.2]
\draw[thick] (0,-1)--(0,0) -- (-2,2);
\draw[thick] (-1,1)--(0,2) ;
\draw[thick] (0,0)--(1,1) ;
\draw [fill] (0,-1) circle [radius=0.035];
\draw [fill] (-2,2) circle [radius=0.035];
\draw [fill] (0,2) circle [radius=0.035];
\draw [fill] (1,1) circle [radius=0.035];
\end{tikzpicture}}}}
\newcommand{\TreeLL}
{\vcenter{\hbox{
\begin{tikzpicture}[yscale=0.2,xscale=0.2]
\draw[thick] (0,-1)--(0,0) -- (-3,3);
\draw[thick] (-1,1)--(0,2) ;
\draw[thick] (-2,2)--(-1,3) ;
\draw[thick] (0,0)--(1,1) ;
\draw [fill] (0,-1) circle [radius=0.035];
\draw [fill] (-2,2) circle [radius=0.035];
\draw [fill] (0,2) circle [radius=0.035];
\draw [fill] (1,1) circle [radius=0.035];
\draw [fill] (-3,3) circle [radius=0.035];
\draw [fill] (-1,3) circle [radius=0.035];
\end{tikzpicture}}}}
\newcommand{\TreeLR}
{\vcenter{\hbox{
\begin{tikzpicture}[yscale=0.2,xscale=0.2]
\draw[thick] (0,-1)--(0,0) -- (-2,2);
\draw[thick] (-1,1)--(1,3) ;
\draw[thick] (0,0)--(1,1) ;
\draw[thick] (0,2)--(-1,3) ;
\draw [fill] (0,-1) circle [radius=0.035];
\draw [fill] (-2,2) circle [radius=0.035];
\draw [fill] (0,2) circle [radius=0.035];
\draw [fill] (1,1) circle [radius=0.035];
\draw [fill] (-1,3) circle [radius=0.035];
\draw [fill] (1,3) circle [radius=0.035];
\end{tikzpicture}}}}
\newcommand{\TreeRR}
{\vcenter{\hbox{
\begin{tikzpicture}[yscale=0.2,xscale=0.2]
\draw[thick] (0,-1)--(0,0) -- (3,3);
\draw[thick] (1,1)--(0,2) ;
\draw[thick] (2,2)--(1,3) ;
\draw[thick] (0,0)--(-1,1) ;
\draw [fill] (0,-1) circle [radius=0.035];
\draw [fill] (2,2) circle [radius=0.035];
\draw [fill] (0,2) circle [radius=0.035];
\draw [fill] (-1,1) circle [radius=0.035];
\draw [fill] (3,3) circle [radius=0.035];
\draw [fill] (1,3) circle [radius=0.035];
\end{tikzpicture}}}}
\newcommand{\TreeRL}
{\vcenter{\hbox{
\begin{tikzpicture}[yscale=0.2,xscale=0.2]
\draw[thick] (0,-1)--(0,0) -- (2,2);
\draw[thick] (1,1)--(-1,3) ;
\draw[thick] (0,0)--(-1,1) ;
\draw[thick] (0,2)--(1,3) ;
\draw [fill] (0,-1) circle [radius=0.035];
\draw [fill] (2,2) circle [radius=0.035];
\draw [fill] (0,2) circle [radius=0.035];
\draw [fill] (-1,1) circle [radius=0.035];
\draw [fill] (1,3) circle [radius=0.035];
\draw [fill] (-1,3) circle [radius=0.035];
\end{tikzpicture}}}}
\newcommand{\TreeCC}
{\vcenter{\hbox{
\begin{tikzpicture}[yscale=0.2,xscale=0.2]
\draw[thick] (0,-1)--(0,0) -- (-2.5,2.5);
\draw[thick] (-1.5,1.5)--(-0.5,2.5) ;
\draw[thick] (1.5,1.5)--(0.5,2.5) ;
\draw[thick] (0,0)--(2.5,2.5) ;
\draw [fill] (0,-1) circle [radius=0.035];
\draw [fill] (-2.5,2.5) circle [radius=0.035];
\draw [fill] (-0.5,2.5) circle [radius=0.035];
\draw [fill] (0.5,2.5) circle [radius=0.035];
\draw [fill] (2.5,2.5) circle [radius=0.035];
\end{tikzpicture}}}}
\newcommand{\TreeAL}
{\vcenter{\hbox{
\begin{tikzpicture}[yscale=0.2,xscale=0.2]
\draw[thick] (0,-1)--(0,0) -- (-2,2);
\draw[thick] (-1,1)--(0,2) ;
\draw[thick] (0,0)--(1,1) ;
\draw[thick] (-1,1)--(-1,2) ;
\draw [fill] (0,-1) circle [radius=0.035];
\draw [fill] (-2,2) circle [radius=0.035];
\draw [fill] (0,2) circle [radius=0.035];
\draw [fill] (1,1) circle [radius=0.035];
\draw [fill] (-1,2) circle [radius=0.035];
\end{tikzpicture}}}}
\newcommand{\TreeRA}
{\vcenter{\hbox{
\begin{tikzpicture}[yscale=0.2,xscale=0.2]
\draw[thick] (0,-1)--(0,0) -- (2,2);
\draw[thick] (1,1)--(0,2) ;
\draw[thick] (0,0)--(-1,1) ;
\draw[thick] (1,1)--(1,2) ;
\draw [fill] (0,-1) circle [radius=0.035];
\draw [fill] (2,2) circle [radius=0.035];
\draw [fill] (0,2) circle [radius=0.035];
\draw [fill] (-1,1) circle [radius=0.035];
\draw [fill] (1,2) circle [radius=0.035];
\end{tikzpicture}}}}
\newcommand{\TreeAR}
{\vcenter{\hbox{
\begin{tikzpicture}[yscale=0.2,xscale=0.2]
\draw[thick] (0,-1)--(0,0) -- (2,2);
\draw[thick] (1,1)--(0,2) ;
\draw[thick] (0,0)--(-1,1) ;
\draw[thick] (0,0)--(0,1) ;
\draw [fill] (0,-1) circle [radius=0.035];
\draw [fill] (2,2) circle [radius=0.035];
\draw [fill] (0,2) circle [radius=0.035];
\draw [fill] (-1,1) circle [radius=0.035];
\draw [fill] (0,1) circle [radius=0.035];
\end{tikzpicture}}}}
\newcommand{\TreeLA}
{\vcenter{\hbox{
\begin{tikzpicture}[yscale=0.2,xscale=0.2]
\draw[thick] (0,-1)--(0,0) -- (-2,2);
\draw[thick] (-1,1)--(0,2) ;
\draw[thick] (0,0)--(1,1) ;
\draw[thick] (0,0)--(0,1) ;
\draw [fill] (0,-1) circle [radius=0.035];
\draw [fill] (-2,2) circle [radius=0.035];
\draw [fill] (0,2) circle [radius=0.035];
\draw [fill] (1,1) circle [radius=0.035];
\draw [fill] (0,1) circle [radius=0.035];
\end{tikzpicture}}}}
\newcommand{\TreeCA}
{\vcenter{\hbox{
\begin{tikzpicture}[yscale=0.2,xscale=0.2]
\draw[thick] (0,-1)--(0,0) -- (-1,1);
\draw[thick] (0,1.5)--(1,2.5) ;
\draw[thick] (0,1.5)--(-1,2.5) ;
\draw[thick] (0,0)--(1,1) ;
\draw[thick] (0,0)--(0,1.5) ;
\draw [fill] (0,-1) circle [radius=0.035];
\draw [fill] (-1,1) circle [radius=0.035];
\draw [fill] (1,2.5) circle [radius=0.035];
\draw [fill] (-1,2.5) circle [radius=0.035];
\draw [fill] (1,1) circle [radius=0.035];
\draw [fill] (0,1) circle [radius=0.035];
\end{tikzpicture}}}}
\newcommand{\TreeC}
{\vcenter{\hbox{
\begin{tikzpicture}[yscale=0.2,xscale=0.2]
\draw[thick] (0,-1.5)--(0,0);
\draw[thick] (0,0)--(1.5,1.5) ;
\draw[thick] (0,0)--(0.5,1.5) ;
\draw[thick] (0,0)--(-0.5,1.5) ;
\draw[thick] (0,0)--(-1.5,1.5) ;
\draw [fill] (0,-1.5) circle [radius=0.035];
\draw [fill] (1.5,1.5) circle [radius=0.035];
\draw [fill] (0.5,1.5) circle [radius=0.035];
\draw [fill] (-1.5,1.5) circle [radius=0.035];
\draw [fill] (-0.5,1.5) circle [radius=0.035];
\end{tikzpicture}}}}
\title{The matching-face formula for the diagonal of Loday associahedra}
\author{Naruki Masuda, Hugh Thomas, Andy Tonks, Bruno Vallette}
\date{doi:10.5802/jep.142; CC BY 4.0}
\begin{document}
\maketitle
\noindent\textit{Editorial note.} The movement lemma1450 prints id:P->P. That literal identity would make every fiber a singleton and every face both lower and upper, so it is not the map used by the ensuing argument. Lines1463–1464 require an edge parallel to e\_i to be vertical, and1474–1475 require line-segment fibers. These uniquely specify the affine-span quotient along span(e\_i), up to affine coordinates in the quotient, with height psi=<x,e\_i>. Retain the printed id and disclose this reading; no new movement proof is inserted.\par\medskip
\noindent\textit{Editorial note.} With the definitions1509, r\_t counts right-leaning leaves minus k and l\_t counts left-leaning leaves minus k. The original contraction recipe1513 gives dim F\_t=r\_t and dim G\_t=l\_t. The displays1509 and1522–1524 exchange these two letters. The recipe, author formulas and proof are retained; the letter convention just stated is the one compatible with the construction and with the inequality r\_s+l\_t<=d used there.\par\medskip
\noindent\textit{Editorial note.} The edge-direction calculation877 has the dot-product difference reversed: the displayed earlier +x/later -x edge pairs as x(v\_i-v\_(i+j))>0. The consecutive-difference sum at888 also has its initial index offset: the source basis886 gives a\_(p+q-1)-a\_p for coordinates p+1,...,p+q-1, with endpoint coefficients0. These are source-determined sign/index readings; the equal-or-zero coefficient argument is unchanged. At1437 the edge is indeed a positive multiple of e\_i, but the scalar is the product of the total weights of the two outside subtrees, not in general the product of two individual leaf weights. The subsequent movement proof uses only positive parallelism. Preserve all printed formulas; do not certify the printed scalar as an exact general value.\par\medskip
\noindent\textit{Editorial note.} At1005 the right-hand map on the lower faces of tilde P is pr, as fixed by the domains and factorization978–982 and the immediately following sentence; tilde-pi has domain Q. This is a cellwise affine homeomorphism on the lower graph, not a claim that the entire piecewise-linear graph is one affine subspace. The original display is retained.\par\medskip
\noindent\textit{Editorial note.} The boundary sentence1651 is compressed. Its order-theoretic input is the classical Tamari face-interval property explicitly listed in standard\_prerequisites. It applies to the s<=t<u vertices and B=G\_t intersect G\_u from the author preceding construction; the conclusion also uses the already proved tightness and full dimension of the current maximal cell. This is an explanation of the scope of the omitted classical convex/order step, not a replacement author proof. The specialized neighbor classification and propagation remain solution obligations.\par\medskip
\noindent\textit{Editorial note.} The counted formula has n>=2 as in582–584; the arity-one operadic unit635 is a separate convention, not an n-2=-1 instance. In product-facet induction omit point factors when choosing a facet and apply the Boolean-fiber facet argument only in positive dimension. The author treats dimensions0 and1 separately at1529; no facet-containing-vertex assertion for a point is made.\par\medskip
\noindent\textit{Editorial scope.} Same counted claim, source spans, original formulas and core obligations as r001. Add six visible source-reading notes and the exact classical Tamari face-interval prerequisite supported by the existing complete C review. The source formulas are not rewritten; no new proof is supplied.\par\medskip
\noindent\textit{Allowed standard prerequisites (editorial).}\begin{itemize}
\item Elementary convex polytope theory, affine projections, face lattices and finite polytopal subdivisions. Standard definitions and elementary face facts only; the specific weighted Loday face decomposition, positivity and tight diagonal construction are retained as solution obligations.
\item Planar binary trees, Tamari order and elementary finite lattice terminology. Definitions and the classical finite Tamari lattice fact; all matching-face exclusion and inclusion arguments must be proved from the retained source.
\item Classical Tamari face-interval property. The binary refinements of a planar/Schroder tree form the Tamari interval between its left-comb and right-comb refinements. This supplies order convexity of a face in the compressed boundary step1651 only. It does not give the diagonal formula, neighbor-face classification, Boolean exclusion or propagation; all those remain the retained author proof.
\end{itemize}\medskip
\section{The approximation of the diagonal of the associahedra}
\subsection{Planar trees, Tamari lattices, and associahedra}

We consider the set $\Tam{n}$ of planar binary (rooted) trees with $n$ leaves, for $n\geq 1$. We read planar binary trees according to gravity, that is from the leaves to the root. The edges of a rooted tree are of three types: the internal edges are bounded by two vertices, the leaves lie at the top and the root at the bottom.

\begin{definition}[Tamari order \cite{Tamari51}]
The \emph{Tamari order} is the partial order, denoted by $<$, on the set of planar binary trees generated by the following covering relation

\[{\vcenter{\hbox{
\begin{tikzpicture}[yscale=0.5,xscale=0.5]
\draw[thick] (0,-1)--(0,0) -- (-2,2);
\draw[thick] (-1,1)--(0,2) ;
\draw[thick] (0,0)--(1,1) ;
\draw [fill] (0,-1) circle [radius=0.014];
\draw [fill] (-2,2) circle [radius=0.014];
\draw [fill] (1,1) circle [radius=0.014];
\draw [fill] (0,2) circle [radius=0.014];

\draw (-2,2) node[above] {$t_1$};
\draw (0,2) node[above] {$t_2$};
\draw (1,1) node[above] {$t_3$};
\draw (0,-1) node[below] {$t_4$};
\end{tikzpicture}}}}
\prec
{\vcenter{\hbox{
\begin{tikzpicture}[yscale=0.5,xscale=0.5]
\draw[thick] (0,-1)--(0,0) -- (2,2);
\draw[thick] (1,1)--(0,2) ;
\draw[thick] (0,0)--(-1,1) ;
\draw [fill] (0,-1) circle [radius=0.014];
\draw [fill] (2,2) circle [radius=0.014];
\draw [fill] (-1,1) circle [radius=0.014];
\draw [fill] (0,2) circle [radius=0.014];

\draw (2,2) node[above] {$t_3$};
\draw (0,2) node[above] {$t_2$};
\draw (-1,1) node[above] {$t_1$};
\draw (0,-1) node[below] {$t_4$};
\end{tikzpicture}}}},
\]
where $t_i$, for $1\leq i\leq 4$, are planar binary trees.
\end{definition}

For every $n\geq 1$, this forms a lattice $(\Tam{n}, <)$. The
\emph{right-leaning} leaves or internal edges are the ones of type $\searrow$ and
the \emph{left-leaning} leaves or internal edges are the ones of type $\swarrow$. So two trees satisfy $s\leq t$ if and only if one goes from $s$ to $t$ by switching
pairs of successive left and right-leaning edges to
corresponding pairs
of successive right and left-leaning edges.

\begin{figure}[htb]
\[\vcenter{\hbox{
\begin{tikzcd}[column sep=0.4cm, row sep=-0.1cm]
& \TreeLL\arrow[rrdd, thin]\arrow[ld, thin]&& \\
\TreeLR\arrow[dd, thin]& && \\
& && \TreeCC \arrow[ldld, thin]\\
\TreeRL\arrow[rd, thin]& && \\
&\TreeRR &&
\end{tikzcd}
}}\]
\caption{The Tamari lattice $(\mathrm{PBT}_{4}, <)$ with minimum at the top.}
\end{figure}

We also consider the set $\PT{n}$ of all planar trees with $n$ leaves, for $n\geq 1$.
Each of these sets forms a lattice (after adjoining a minimum) under the following partial order: a planar tree $s$ is less than a planar tree $t$, denoted $s\subset t$,
if $t$ can be obtained from $s$ by a sequence of edge contractions.

\begin{definition}[Associahedra]
For any $n\geq 2$, an \emph{$(n-2)$-dimensional associahedron} is a polytope whose face lattice is isomorphic to the lattice of planar trees with $n$ leaves.
\begin{figure}[htb]
\[
\begin{tikzpicture}[scale=0.3]
\draw[thick] (-5,0)--(0,4)--(5,0)--(2.5,-4)--(-2.5,-4)--cycle;
\draw[fill, opacity=0.12] (-5,0)--(0,4)--(5,0)--(2.5,-4)--(-2.5,-4)--cycle;
\draw (-5,0) node[left] {$\TreeLL$};
\draw (5,0) node[right] {$\TreeRR$};
\draw (0,4) node[above] {$\TreeCC$};
\draw (-2.5,-4) node[below left] {$\TreeLR$};
\draw (2.5,-4) node[below right] {$\TreeRL$};
\draw (-3.8,-2.6) node[left] {$\TreeAL$};
\draw (3.7,-2.6) node[right] {$\TreeRA$};
\draw (-2.3,1.8) node[above left] {$\TreeLA$};
\draw (2.3,1.8) node[above right] {$\TreeAR$};
\draw (0,-4.2) node[below] {$\TreeCA$};
\draw (0,-0.4) node {$\TreeC$};
\end{tikzpicture}
\]
\caption{A 2-dimensional associahedron.}
\end{figure}
\end{definition}

The codimension of a face is equal to the number of internal edges of the corresponding planar tree.
The $1$-skeleton of an associahedron gives the Hasse diagram of the Tamari lattice.

The operation of grafting a tree $t$ at the $i^\text{th}$-leaf of a tree $s$ is denoted by $s\circ_i t$.
These maps endow the collections of planar (binary) trees with a non-symmetric operad structure.
We denote the corolla with $n$ leaves by $c_n$, \ie the tree with one vertex and no internal edge. The facets of an $(n-2)$-dimensional associahedron are labeled by
the two-vertex planar trees $c_{p+1+r} \circ_{p+1} c_q$, for $p+q+r=n$ with $2\leq q\leq n-1$.

\subsection{Loday realizations of the associahedra}

An example of realization of the associahedra can be given as follows; it is a weighted generalization of the one given by Jean-Louis Loday in \cite{Loday04a}. Notice that Loday realizations produced as convex hulls are the same polytopes as Shnider--Stasheff \cite{StasheffShnider97} produced by intersections of half-spaces, or equivalently by truncations of standard simplices.

\begin{definition}[Weighted planar binary tree]
A \emph{weighted planar binary tree} is a pair $(t, \omega)$ made up of a planar binary tree $t\in \Tam{n}$ with $n$ leaves having some weight $\omega= (\omega_1, \ldots, \omega_n) \in \ZP^n$. We call
$\omega$ the \emph{weight} and
$n$ the \emph{arity} of the tree $t$ or the \emph{length} of the weight $\omega$.
\end{definition}

Let $(t, \omega)$ be a weighted planar binary tree with $n$ leaves. We order its $n-1$ vertices from left to right. At the $i^\text{th}$ vertex, we consider the sum $\alpha_i$ of the weights of the leaves supported by its left input and
the sum $\beta_i$ of the weights of the leaves supported by its right input.
Multiplying these two numbers, we consider the associated string which defines the following point:\vspace*{-3pt}
\[M(t, \omega) \coloneqq \big(\alpha_1\beta_1, \alpha_2\beta_2, \ldots, \alpha_{n-1}\beta_{n-1}\big)\in \RR^{n-1}. \]

\begin{definition}[Loday realization]
The \emph{Loday realization of weight $\omega$} is the polytope\vspace*{-3pt}
\[\K_\omega \coloneqq \conv \big\{M(t, \omega)\mid t\in \Tam{n} \big\}\subset \RR^{n-1}.\]
\end{definition}

The Loday realization associated to the standard weight $(1, \ldots, 1)$ is simply denoted by $\K_n$.
By convention, we define the polytope $K_\omega$, with weight $\omega=(\omega_1)$ of length $1$, to be made up of one point labeled by the trivial tree $|$.
\begin{figure}[htb]
\[
\begin{tikzpicture}[scale=0.5, JLL3d]
\draw[->] (0,4,-2.5)--(1.5,4,-2.5) node[right] {$\vec e_1$};
\draw[->] (0,4,-2.5)--(0,5,-2.5) node[left] {$\vec e_2$};
\draw[->] (0,4,-2.5)--(0,4,-1.5) node[above] {$\vec e_3$};

\draw[thick] (3,4,3)--(3,0,3)--(3,0,-1)--(3,2,-1)--(3,4,1)--cycle;
\draw[thick] (-3,4,3)--(-3,-4,3)--(-1,-4,3);
\draw[thick] (-1,-4,3)--(3,0,3);
\draw[thick] (-3,-4,3)--(-1,-4,3)--(-1,-4,-3)--(1,-2,-3);
\draw[thick] (1,-2,-3)--(3,0,-1);
\draw[thick] (-1,-4,-3)--(-1,-4,3);
\draw[thick] (-1,-4,-3)--(1,-2,-3)--(1,0,-3)--(-3,0,-3);
\draw[thick, opacity=0.2] (-3,4,1)--(-3,0,-3)--(-3,-4,-3)--(-3,-4,3)--(-3,4,3);
\draw[thick, opacity=0.2] (-3,-4,-3)--(-1,-4,-3);
\draw[thick] (-3,4,3)--(3,4,3)--(3,4,1);
\draw[thick] (-3,-4,3)--(-3,4,3)--(-3,4,1)--(3,4,1);
\draw[thick] (-3,4,3)--(-3,4,1)--(-3,0,-3);
\draw[thick] (1,0,-3)--(3,2,-1);

\draw[fill, opacity=0.18] (-3,4,1)--(-3,4,3)--(3,4,3)--(3,4,1)--cycle;
\draw[fill, opacity=0.18] (-1,-4,3)--(3,0,3)--(3,4,3)--(-3,4,3)--(-3,-4,3)--cycle;
\draw[fill, opacity=0.18] (-3,0,-3)--(-3,4,1)--(3,4,1)--(3,2,-1)--(1,0,-3)--cycle;
\draw[fill, opacity=0.06] (1,-2,-3)--(3,0,-1)--(3,0,3)--(-1,-4,3)--(-1,-4,-3)--cycle;
\draw[fill, opacity=0.06] (1,-2,-3)--(1,0,-3)--(3,2,-1)--(3,0,-1)--cycle;
\draw[fill, opacity=0.12] (3,4,3)--(3,0,3)--(3,0,-1)--(3,2,-1)--(3,4,1)--cycle;
\end{tikzpicture}
\]
\caption{The Loday realization $\K_5$, see \cref{prop:OrientationVector} for the definition of the $\vec e_i$.}
\end{figure}

In the sequel, we will need the following key properties of these polytopes. They mainly come from \cite{Loday04a} and \cite[\S6]{Forcey08}.

\begin{proposition}\label{prop:PropertiesKLoday}
The Loday realization $\K_\omega$ satisfies the following properties.
\begin{enumerate}
\item\label{prop:PropertiesKLoday1} It is contained in the hyperplane with equation\vspace*{-3pt}
\[
\sum_{i=1}^{n-1} x_i = \sum_{1\leq k< l\leq n} \omega_k \omega_l.
\]

\item\label{prop:PropertiesKLoday2} Let $p+q+r=n$ with $2\leq q\leq n-1$. For any $t\in \Tam{n}$, the point $M(t, \omega)$ is contained in the half-space defined by the inequality\vspace*{-3pt}
\[
x_{p+1}+\cdots+x_{p+q-1}\geq \sum_{p+1\leq k<l\leq p+q} \omega_k \omega_l,
\]
with equality if and only if the tree $t$ can be decomposed as $t=u\circ_{p+1} v$, where $u\in\Tam{p+1+r}$ and $v\in \Tam{q}$.

\item\label{prop:PropertiesKLoday3} The polytope $K_\omega$ is the intersection of the hyperplane of \eqref{prop:PropertiesKLoday1} and the half-spaces of \eqref{prop:PropertiesKLoday2}.

\item\label{prop:PropertiesKLoday4} The face lattice $(\La(\K_\omega), \subset)$ is equal to the lattice $(\mathrm{PT}_n, \subset)$
of planar trees with $n$ leaves.

\item\label{prop:PropertiesKLoday5} Any face of a Loday realization is isomorphic to a product of Loday realizations, via a permutation of coordinates.
\end{enumerate}
\end{proposition}

\begin{proof}
This proposition is a weighted version of the results of \cite{Loday04a}, except for Point~\eqref{prop:PropertiesKLoday5}, which actually prompted the introduction of this more general notion.
\leavevmode
\begin{enumerate}
\item This is straightforward from the definition.
\item This is straightforward too.
\item Let us denote by $P$ the polytope defined by the intersection of the hyperplane of \eqref{prop:PropertiesKLoday1} and the half-spaces of \eqref{prop:PropertiesKLoday2}.
One can see that the points $M(t, \omega)$, for $t\in \Tam{n}$, are vertices of $P$, since they are defined by a system of $n-1$ independent linear equations: the one of type \eqref{prop:PropertiesKLoday1} and $n-2$ of type \eqref{prop:PropertiesKLoday2}.
In the other way round, any vertex of $P$ is characterized by a system $n-1$ independent linear equations with the one of type \eqref{prop:PropertiesKLoday1} and $n-2$ of type \eqref{prop:PropertiesKLoday2}. We claim that any pair of equations
\begin{align*}
x_{p+1}+\cdots+x_{p+q-1}&= \sum_{p+1\leq k<l\leq p+q} \omega_k \omega_l\\
\tag*{and}
x_{p'+1}+\cdots+x_{p'+q'-1}&= \sum_{p'+1\leq k<l\leq p'+q'} \omega_k \omega_l
\end{align*}
of type \eqref{prop:PropertiesKLoday2} appearing here is such that the corresponding intervals
\[
[p+1, \ldots, p+q-1]\quad\text{and}\quad [p'+1, \ldots, p'+q'-1]
\]
are either nested or disjoint.
If not, we are in the configuration:
\[
p+1<p'+1\leq p+q+1<p'+q'+1.
\]
Using these equalities and the defining inequalities of $P$, one can get
\begin{align*}
x_{p'+1}+\cdots+x_{p+q-1}&\leq
\sum_{p+1\leq k<l\leq p+q}\hspace*{-3mm} \omega_k \omega_l
+\sum_{p'+1\leq k<l\leq p'+q'}\hspace*{-3mm} \omega_k \omega_l
-\sum_{p+1\leq k<l\leq p'+q'}\hspace*{-3mm} \omega_k \omega_l \\
&=
\sum_{p'+1\leq k<l\leq p+q}\hspace*{-3mm} \omega_k \omega_l
-
\sum_{\substack{p+1\leq k\leq p'\\ p+q+1\leq l\leq p'+q'}}\hspace*{-3mm} \omega_k \omega_l
< \sum_{p'+1\leq k<l\leq p+q}\hspace*{-3mm} \omega_k \omega_l,
\end{align*}
which contradicts the definition of $P$.
The solution of a system of equations as above, where the defining intervals are nested or disjoint, is a point
$M(t,\omega)$, with
$t\in \Tam{n}$.

\item Point~\eqref{prop:PropertiesKLoday2} shows that the facets of $\K_\omega$ correspond bijectively to two-vertex planar trees: the facet labeled by
$c_{p+1+r}\circ_{p+1} c_q$, for $p+q+r=n$ and $2\leq q \leq n-1$, is the convex hull of the points $M(t, \omega)$ associated to planar binary trees of the form $t=u\circ_{p+1} v$, for $u\in\Tam{p+1+r}$ and $v\in \Tam{q}$.
Any face of $\K_\omega$ of codimension $k$, for
$0\leq k \leq n-2$, is defined as an intersection of $k$ facets. The above description of facets gives that the set of faces of codimension $k$ is bijectively labeled by planar trees with~$k$ internal edges: the face corresponding to such a planar tree $t$ is the convex hull of points $M(s, \omega)$, for $s\subset t$. With the top dimensional face labeled by the corolla $c_n$, the statement is proved.

\item The proof of the above point shows that it is enough to treat the case of the facets.
Let $p+q+r=n$ with $2\leq q \leq n-1$. We consider the following two weights
\[
\overline{\omega}\coloneqq (\omega_1, \ldots, \omega_{p}, \omega_{p+1}+\cdots+\omega_{p+q}, \omega_{p+q+1}, \ldots, \omega_{n})
\quad \text{and} \quad
\widetilde{\omega}\coloneqq (\omega_{p+1}, \ldots, \omega_{p+q}).
\]
The image of $\K_{\overline{\omega}}\times \K_{\widetilde{\omega}}\hto \K_\omega$ under the isomorphism
\begin{align*}
\begin{array}{rccc}
\Theta\ : & \RR^{p+r}\times \RR^{q-1} &\xrightarrow{~\cong~} &\RR^{n-1}\\
&(x_1, \ldots, x_{p+r})\times (y_1, \ldots, y_{q-1}) & \longmapsto&
(x_1, \ldots, x_{p}, y_1, \ldots, y_{q-1}, x_{p+1}, \ldots, x_{p+r})
\end{array}
\end{align*}
is equal to the facet labeled by the planar tree $c_{p+1+r}\circ_{p+1} c_q$.\qedhere
\end{enumerate}
\end{proof}

\subsection{The category of polytopal subdivisions}
The proposed notions of category of polytopes present in the literature only allow affine maps, which is too restrictive for our purpose: the facets of the Loday realizations of associahedra with standard weights are not affinely equivalent to the product of lower realizations with standard weights. In order to introduce a more suitable category, we begin with the following definition, which extends the classical notion of simplicial complex.

\begin{definition}[Polytopal complex]
A \emph{polytopal complex} is a finite collection $\sC$ of polytopes of $\RR^n$ satisfying the following conditions:
\begin{enumerate}
\item the empty polytope $\emptyset$ is contained in $\sC$,
\item $P \in \sC$ implies $\La(P) \subset \sC$,
\item $P, Q\in \sC$ implies $P\cap Q\in \La(P)\cap \La(Q)$.
\end{enumerate}
\end{definition}
Any polytope $P$ gives an example of polytopal complex $\La(P)$ made up of all its faces.
A subcomplex of a polytopal complex $\sC$ is a subcollection
$\mathcal{D}\subset \sC$ which forms a polytopal complex.
The \emph{underlying set} of a collection $\sC$ is given by the union $|\mathcal{C}|\coloneqq\bigcup_{P\in \sC} P\subset \RR^n$.

\begin{definition}[Polytopal subdivision]
A \emph{polytopal subdivision} of a polytope $P$ is a polytopal complex $\sC$
whose underlying set $|\sC|$ is equal to $P$.
\end{definition}

The face poset $\La(\sC)$ of a polytopal complex is defined in the obvious way.
We say that two polytopal complexes are \emph{combinatorially equivalent} when their face posets are isomorphic.
Now let us introduce the category we will work in.

\begin{definition}[The category $\PolySub$]
The category $\PolySub$ is made up of the following data.
\begin{itemize}
\item\emph{Objects}: An object is a $d$-dimensional polytope $P$ in the $n$-dimensional Euclidian space $\RR^n$, for any $0\leq d\leq n$.
\item\emph{Morphisms}: A continuous map $f: P\to Q$ is a morphism when
it sends $P$ homeomorphically to the underlying set $|\mathcal{D}|$ of a polytopal subcomplex $\mathcal{D}\subset \La(Q)$ of~$Q$
such that $f^{-1}(\mathcal D)$ defines a polytopal subdivision of $P$.
\end{itemize}
\end{definition}

There exists obvious forgetful functors from the category $\PolySub$ to the category $\Top$ of topological spaces with continuous maps and to the category $\CW$ of CW complexes with cellular maps. The latter functor is well-defined since any morphism in $\PolySub$ is automatically cellular.
An isomorphism $P\cong Q$ in this category is a cell-respecting homeomorphism which induces
a combinatorial equivalence $\La (P)\cong \La(Q)$.

\begin{lemma}
The category $\PolySub$ endowed with the direct product $\times$ and the zero-dimensional polytope made up of one point is a symmetric monoidal category.
\end{lemma}

\begin{proof}
The verification of the axioms is straightforward.
\end{proof}

This extra structure allows one to consider operads in the category $\PolySub$.
Since the cellular chain functor $\PolySub \to \textsf{dgMod}_{\ZZ}$ is
strong symmetric monoidal, it induces a functor from the category of polytopal (non-symmetric) operads to the category of differential graded (non-symmetric) operads.\enlargethispage{\baselineskip}%

\section{Canonical diagonal for positively oriented polytopes}\label{Sec:CanDia}

\Subsection{Positively oriented polytopes}

\skpt
\begin{definition}[Positively oriented polytope]
\begin{enumerate}
\item An \emph{oriented polytope} is a polytope $P\subset \RR^n$ endowed with a vector $\vec v\in \RR^n$ such that no edge of $P$ is perpendicular to $\vec v$, see \cref{Fig:PosOr}.
\begin{figure}[htb]
\[\vcenter{\hbox{
\begin{tikzpicture}[xscale=0.5, yscale=0.4]
\draw[thick] (0,3)--(-2,1)--(-2,-1)--(0,-3)--(3,0)--cycle;
\draw[->] (4.5,1.5)--(4.5,-1.5);
\draw (4.5,0) node[right] {$\vec v$};
\end{tikzpicture}
}}\]
\caption{An oriented polytope.}\label{Fig:PosOr}
\end{figure}

\item A \emph{positively oriented polytope} is an oriented polytope $(P, \vec v)$ such that the intersection polytope
\[
\big(P\cap \rho_z P, \vec v\big)
\]
is oriented, where $\rho_z\coloneqq 2z-P$ stands for the reflection with respect to any point $z\in P$, see \cref{Fig:Diag}.
\end{enumerate}
\end{definition}

The data of an orientation vector induces a poset structure on the set of vertices $\Ve(P)$ of $P$, which is the transitive closure of the relation induced by the oriented edges of the 1-skeleton.

\begin{definition}[Well-oriented realization of the associahedron]
A \emph{well-oriented reali\-zation of the associahedron} is a positively oriented polytope which realizes the associahedron and such that the orientation vector induces the Tamari lattice on the set of vertices.
\end{definition}

\begin{proposition}\label{prop:OrientationVector}
Let $\omega$ be a weight of length $n$.
Any vector $\vec v=(v_1, \ldots, v_{n-1})\in \RR^{n-1}$ satisfying $v_1>v_2>\cdots >v_{n-1}$ induces a
well-oriented realization $(\K_\omega, \vec v)$ of the associahedron.
\end{proposition}

\begin{proof}
Let us first prove that such an orientation vector $\vec v$ induces the Tamari lattice.
Let $s\prec t$ in the Tamari lattice. The corresponding edge in $\K_\omega$ is of the form
\[\overrightarrow{M(s,\omega)M(t,\omega)}=(0, \ldots, 0, x, 0, \ldots, 0, -x, 0, \ldots, 0), \]
with $x>0$, which implies
$\big\langle\overrightarrow{M(s,\omega)M(t,\omega)}, \vec v\big\rangle= x(v_{j+i}-v_i)>0$.
This also proves that $(K_\omega,\vec v)$ is oriented.

Let us now prove that $(\K_\omega,\vec v)$ is positively oriented.
We denote by $\vec n\coloneqq (1, \ldots, 1)$ and $\vec n_F\coloneqq (0, \ldots, 0,\allowbreak {1}, \ldots, {1}, 0, \ldots, 0)$ the normal vectors of a facet $F$, given by Points~\eqref{prop:PropertiesKLoday1} and \eqref{prop:PropertiesKLoday2} of \cref{prop:PropertiesKLoday}.
Since edges of $\K_\omega\cap \rho_z \K_\omega$
are one-dimensional intersections of facets of
$\K_\omega$ or $\rho_z \K_\omega$, their directions $\vec d$ are the unique solutions to a system of equations of type $\langle \vec n, \vec d\rangle=0$ and $\langle \vec n_F, \vec d\rangle=0$, for some set of facets $F$.
The first equation imposes $\vec d=\sum_{j=1}^{n-2} a_j \vec e_j$, where
$\vec e_j\coloneqq (0, \ldots, 1, -1, \ldots, 0)$, in which $1$ is in the $j$-th place for $1\leq j\leq n-2$.
The second equation is equivalent to one of the following three constraints
$a_{p+1}=0$, $a_{p+q-1}=0$, or $a_{p+q-1}=a_{p+1}$.
Therefore, $\vec d$ is of the form $\lambda\left(\vec e_{j_1}+\cdots +\vec e_{j_k}\right)$, with $\lambda \in \RR\setminus \{0\}$, and so
$\ang{\vec d, \vec v}\neq 0$.
\end{proof}

\subsection{Construction of diagonal maps}

Before getting into specific argument on the associahedra, we construct a diagonal map $\triangle: P\to P\times P$ for any positively oriented polytope $(P, \vec v)$. Let $\tp_{\vec v} P$ (\resp $\bm_{\vec v} P$) denote the top (\resp bottom) vertex of $P$ with respect to the orientation vector $\vec v$.

\begin{definition}[Diagonal of a positively oriented polytope]\label{def:Diag}
The \emph{diagonal} of a positively oriented polytope $(P, \vec v)$ is defined by
\begin{align*}
\begin{array}{rlcl}
\triangle_{(P,\vec v)}\ : & P &\dpl\to &P\times P\\
&z & \dpl\mto&
\bigl(\bm_{\vec v}(P\cap \rho_z P),\, \tp_{\vec v}(P\cap \rho_z P)\bigr).
\end{array}
\end{align*}
\begin{figure}[htb]
\[
\begin{tikzpicture}[xscale=0.6, yscale=0.5]
\draw[thick] (1,3)--(-1,1)--(-1,-1)--(1,-3)--(4,0)--cycle;
\draw[thick] (-1,3)--(1,1)--(1,-1)--(-1,-3)--(-4,0)--cycle;

\draw[thick] (0,-2)--(-1,-1)--(-1, 1)--(0, 2)--(1, 1)--(1,-1)--cycle;
\draw[fill, opacity=0.15] (0,-2)--(-1,-1)--(-1, 1)--(0, 2)--(1, 1)--(1,-1)--cycle;

\draw[->] (5,1.2)--(5,-1.2);
\draw (5,0) node[right] {$\vec v$};
\draw (0,0) node {$\bullet$};
\draw (0,2) node {$\bullet$};
\draw (0,-2) node {$\bullet$};

\draw (0,0) node[right] {$z$};
\draw (0,1.9) node[below] {$x$};
\draw (0,-1.9) node[above] {$y$};

\draw (3,-2) node {$P$};
\draw (-3,-2) node {$\rho_z P$};
\end{tikzpicture}
\]
\caption{The diagonal map
$$\hspace*{-3cm}\triangle_{(P,\vec v)}(z)=\bigl(\bm_{\vec v}(P\cap \rho_z P),\, \tp_{\vec v}(P\cap \rho_z P)\bigr)=(x,y).$$}\label{Fig:Diag}
\end{figure}
\vspace*{-\baselineskip}
\end{definition}

Let $\beta$ be the middle-point map $P\times P\to P; (x,y)\mto \frac{x+y}{2}$.
With the notation $\pr_i$ for the projection to the $i$-th coordinate, we have $\pr_1\beta^{-1}(z)=\pr_2\beta^{-1}(z) = P\cap \rho_z P$.
The diagonal $\triangle_{(P,\vec v)}$ is a section of $\beta$ since $P\cap \rho_z P$ is symmetric with respect to the point $z$.

Since the Loday realizations $K_\omega$ of the associahedra are positively oriented when the orientation vector $\vec v$ has decreasing coordinates, the above formula endows them with diagonal maps, which do not depend on the choice of such an orientation vector.

\begin{proposition}\label{prop:IndepVect}
Let $\vec v$ and $\vec w$ be two vectors of $\RR^{n-1}$ with decreasing coordinates, then \[\triangle_{(K_\omega, \vec v)}=\triangle_{(K_\omega, \vec w)},\]
for every weight $\omega$ of length $n$.
\end{proposition}

\begin{proof}
The argument given in the proof of \cref{prop:OrientationVector} shows that the formula
\[\bigl(\bm_{\vec v}(\K_\omega\cap \rho_z \K_\omega),\, \tp_{\vec v}(\K_\omega\cap \rho_z \K_\omega)\bigr)=
\bigl(\bm_{\vec w}(\K_\omega\cap \rho_z \K_\omega),\, \tp_{\vec w}(\K_\omega\cap \rho_z \K_\omega)\bigr)
\]
produces the same pair for any orientation vector with decreasing coordinates.
\end{proof}

We denote by $\triangle_\omega : \K_\omega\to \K_\omega \times \K_\omega$ the diagonal map given by any such orientation vector and we use the simple notation $\tp$ and $\bm$ for associahedra.
In the case of the standard weight $\omega=(1, \ldots, 1)$, we denote the diagonal simply by $\triangle_n : \K_n \to \K_n\times \K_n$.

\begin{lemma}\label{Lemma:DiagOnFace}
Any face $F$ of a positively oriented polytope $(P,\vec v)$ is positively oriented once equipped with the orientation vector $\vec v$. Moreover, the two diagonals agree: $\triangle_{(P, \vec v)}(z)=\triangle_{(F, \vec v)}(z)$, for any $z\in F$.
\end{lemma}

\begin{proof}
This follows from the relation
$P\cap \rho_z P= F\cap \rho_z F$, for any $z\in F$.
\end{proof}

\subsection{Polytopal subdivision induced by the diagonal}\label{subsec:PolySubDiv}
The above formula for the diagonal $\triangle$ actually defines a morphism in the category $\PolySub$.
To prove this property, we use some ideas coming from the theory of fiber polytopes \cite{BilleraSturmfels92}, see also \cite[Chap.\,9]{Ziegler95}.

Let $\pi: Q\twoheadrightarrow P$ be an affine projection of polytopes with
$P\subset \RR^p$ and $Q\subset \RR^q$.
We~denote by $Q_z\coloneqq Q\cap \pi^{-1}(z)$ the fiber above $z\in P$.
To any linear form $\psi\in (\RR^q)^{\ast}$, we~associate a
collection
$\sF^\psi\subset \La(Q)$ as follows.
We first factor the affine projection~as\vspace*{-3pt}
\begin{align*}
\pi&=\mathrm{pr}\circ \tilde{\pi},& \text{where} \quad
\tilde{\pi}&\coloneqq (\pi, \psi): Q\twoheadrightarrow \tilde{P}\\
\tag*{and}
\mathrm{pr}&: \tilde{P}\twoheadrightarrow P;\, (z, t)\mto z,&
\text{where}\quad \tilde{P}&\coloneqq \big\{(\pi(x), \psi(x))\mid x\in Q\big\}\subset \RR^{p+1}.
\end{align*}
Next let $\La^{\downarrow}(\tilde P) \subset \La\big(\tilde{P}\big)$ be the subcomplex of lower faces, \ie $F\in \La^{\downarrow}(\tilde P)$ if and only if any $(z, t)\in F$ satisfies the equation
$t = \min \psi(Q_z)$. Since the preimage of any face by a projection of polytopes is again a face, this defines a collection\vspace*{-3pt}\enlargethispage{.5\baselineskip}%
\[\sF^\psi \coloneqq
\big\{Q\cap \tilde{\pi}^{-1}(F) \mid F\in \La^{\downarrow}(\tilde P)
\big\}
\subset \La(Q).\]
(This is in general not a polytopal complex since it is not stable under faces.)
In the case of the point $P=\{\ast\}$, the unique face contained in $\sF^\psi$ is by definition\vspace*{-3pt}
\[Q^\psi\coloneqq\{x\in Q \mid \psi (x) = \min \psi(Q)\}.\]

\begin{proposition}\label{prop:CoherentSubdivision}
The collection $\sF^\psi \subset \La(Q)$ satisfies the following properties.
\begin{enumerate}
\item\label{prop:CoherentSubdivision1} The polytopal complex $\pi(\sF^\psi)\coloneqq\left\{\pi(F)\mid F\in \sF^\psi \right\}$ is a polytopal subdivision of $P$.
\item\label{prop:CoherentSubdivision2} For any $z\in P$, the fiber satisfies $(\sF^\psi)_z \coloneqq\pi^{-1}(z) \cap\left| \sF^\psi\right|=\left(Q_z\right)^\psi$.
\end{enumerate}
\end{proposition}

\skpt
\begin{proof}
\begin{enumerate}
\item By definition $\pi\left(\sF^\psi\right) = \tilde\pi\big(\La^{\downarrow}(\tilde P)\big)$. The right-hand side defines a polytopal subdivision of $P$, since the restriction of the map $\pr$ is a linear homeomorphism from a polytopal complex.
\item This is clear from the definition.\qedhere
\end{enumerate}
\end{proof}

\skpt
\begin{definition}[Coherent and tight subdivisions]\label{def:Coherent}
\begin{enumerate}
\item A subcollection $\sF\subset\La(Q)$ is called a \emph{coherent subdivision of $P$} when it is of the form $\sF^\psi$ for some $\psi\in (\RR^q)^{\ast}$.
\item A coherent subdivision $\sF$ is called \emph{tight} when, the faces $F$ and $\pi(F)$ have the same dimension, for every $F\in \sF$.
\end{enumerate}
\end{definition}

A coherent subdivision $\sF$ is tight if and only if, for any $z\in P$, the fiber $\sF_z = (Q_z)^\psi$ is a point, by Point~\eqref{prop:CoherentSubdivision2} of \cref{prop:CoherentSubdivision}. (This is also equivalent to $\sF$ being a polytopal complex.)
Therefore a tight coherent subdivision can be identified with the unique piecewise-linear section of $\pi|_Q$ which minimizes $\psi$ in each fiber.
By the Point~\eqref{prop:CoherentSubdivision1} of \cref{prop:CoherentSubdivision}, this section is a morphism of the category $\PolySub$.

We apply these results to the projection\vspace*{-3pt}
\[\beta : P\times P\to P ;\quad (x, y)\mto \frac{x+y}{2}.\]
and to the linear form $\psi(x, y)\coloneqq \ang{x-y, \vec v}$ in order to obtain the following proposition.

\begin{proposition}\label{prop:DiaginPoly}
If $(P, \vec v)$ is a positively oriented polytope, the diagonal map $\triangle_{(P, \vec v)}: P\to P\times P$ is a morphism in the category $\PolySub$.
\end{proposition}

\begin{proof}
For any $z\in P$, the fiber $\beta^{-1}(z)$ is the set of pairs $(x, y)\in P\times P$ such that $x+y = 2z$. Since the sum of $x$ and $y$ is constant, $\psi(x, y)$ is minimized in the fiber when $\ang{x, \vec v}$ is minimized, or equivalently, $\ang{y, \vec v}$ is maximized. On both coordinates, $\beta^{-1}(z)$ projects down to the intersection $P\cap (2z-P)$, which is oriented by definition.
So the fiber $\beta^{-1}(z)$ admits a unique minimal element $\left(\bm_{\vec v}(P\cap \rho_z P),\, \tp_{\vec v}(P\cap \rho_z P)\right)$ with respect to $\psi$. In the end, the section defined by the tight coherent subdivision agrees with the definition of the diagonal map $\triangle_{(P, \vec v)}$ given in \cref{def:Diag}.
\end{proof}

\begin{corollary}\label{cor:DIMENSION}
The image of $\triangle_{(P,\vec v)}$ is contained in $\sk_n (P\times P)$, where $n$ is the dimension of $P$. In particular, if one of two components of $\triangle (z)$ lies in the interior of $P$, then the other component is either $\tp_{\vec v} P$ or $\bm_{\vec v} P$.
\end{corollary}

\begin{remark}
Notice that the diagonal map $\triangle_{(P,\vec v)}$ is fiber-homotopic to the usual diagonal $x\mto (x, x)$.
\end{remark}

We denote by $\sF_{(P, \vec v)}$ the corresponding tight coherent subdivision of $P$ and by $\beta(\sF_{(P, \vec v)})$ the polytopal subdivision of $P$.
In the case of the Loday realizations of the associahedra, \cref{prop:IndepVect} shows that they do not depend on the orientation vector (with decreasing coordinates), so we use the simple notations $\sF_{\omega}$ and $\beta(\sF_\omega)$.

\section{The magical formula}\label{sec:MagicalFormula}

\begin{theorem}[Magical formula]\label{thm:MagicFormula}
For any Loday realization of the associahedra, the approximation of the diagonal satisfies
\[{\im \triangle_n
=\bigcup_{\substack{\tp F\leq \bm G\\ \dim F+\dim G=n-2}} F\times G}.\]
\begin{figure}[htb]
\[
\begin{tikzpicture}[scale=0.8]
\draw[thick] (-5,0)--(0,4)--(5,0)--(2.5,-4)--(-2.5,-4)--cycle;
\draw[fill, opacity=0.12] (-5,0)--(0,4)--(5,0)--(2.5,-4)--(-2.5,-4)--cycle;

\draw (-2.5,2)--(0,0)--(-1.5,-2)--(-3.75,-2);
\draw (2.5,2)--(0,0)--(1.5,-2)--(3.75,-2);
\draw (-1.5,-2)--(0,-4)--(1.5,-2);

\draw (-2.7,-0.25) node {$\TreeLL\ \times\TreeC$};
\draw (2.7,-0.25) node {$\TreeC\ \times\TreeRR$};
\draw (0,-2.2) node {$\TreeAL\ \times\TreeRA$};
\draw (0,2) node {$\TreeLA\ \times\TreeAR$};
\draw (-2.25,-3) node {$\TreeAL\ \times\TreeCA$};
\draw (2.15,-3) node {$\TreeCA\ \times\TreeRA$};

\draw (-5,0) node[left] {$\TreeLL$};
\draw (5,0) node[right] {$\TreeRR$};
\draw (0,4) node[above] {$\TreeCC$};
\draw (-2.5,-4) node[below left] {$\TreeLR$};
\draw (2.5,-4) node[below right] {$\TreeRL$};
\draw (-4,-2) node[left] {$\TreeAL$};
\draw (4,-2) node[right] {$\TreeRA$};
\draw (-2.5,2) node[above left] {$\TreeLA$};
\draw (2.5,2) node[above right] {$\TreeAR$};
\draw (0,-4) node[below] {$\TreeCA$};
\end{tikzpicture}
\]
\caption{The polytopal subdivision $\beta\left(\sF_4\right)$ of $\K_4$.}
\end{figure}
\end{theorem}

\subsection{First step: $\im \triangle_n \subset \bigcup F\times G$}\label{Subsec:Subset}
We prove this property more generally for any \emph{product} $P \coloneqq K_{\omega_1}\times \cdots \times K_{\omega_k}$ of Loday realizations of the associahedra. Recall that
$P\subset \RR^{N}$, where $N\coloneqq n_1+\cdots+n_k-k$, and $d\coloneqq \dim P = n_1+\cdots+n_k-2k$.
The set~$\Ve(P)$ of vertices of $P$ coincides with
$\Tam{n_1}\times \cdots \times \Tam{n_k}$.
By \cref{prop:OrientationVector}, any vector
$\vec v= (v_1, v_2, \ldots, v_N)$ satisfying
\[v_1>\cdots>v_{n_1-1},\ \ v_{n_1}>\cdots>v_{n_1+n_2-2},\ldots, \quad v_{n_1+\cdots+n_{k-1}-k+2}>\cdots>v_{N}\ \]
makes $(P, \vec v)$ into a positively oriented polytope with $1$-skeleton isomorphic to the product of Tamari lattices.

We consider the map $L=(L_i)_{1\leq i\leq m-2} : \Tam{m}\to \Bool{m-2}$
to the Boolean lattice
defined by
\[
L_i(t)\coloneqq\begin{cases}
0& \text{if the $(i+1)$-th leaf is left-leaning $\swarrow$,}\\
1& \text{if the $(i+1)$-th leaf is right-leaning $\searrow$}\,.
\end{cases}
\]
We extend it to a map $L=(L_i)_{1\leq i \leq d}: \Ve(P)\to \Bool{d}$. We consider the collection of vectors
$\vec e_1, \ldots, \vec e_d$ in $\RR^N$ defined by
$\vec e_i\coloneqq (0, \ldots, 1, -1, \ldots, 0)$, where $1$ is in the
place $a_i$, with $\{a_i\}_{1\leq i \leq d}$ the increasing sequence found by starting with $1,\dots, N$ and deleting the values $n_1+\cdots+n_j-j$, for $j = 1,\ldots, k$.

\begin{lemma}\label{lemma:propertiesofL}
The following properties are satisfied:
\begin{enumerate}
\item\label{lemma:propertiesofL1} $s\leq t$ $\Rightarrow L(s)\leq L(t)$,

\item\label{lemma:propertiesofL2} $\ang{\vec e_i, \vec v} >0$, for $1\leq i \leq d$,

\item\label{lemma:propertiesofL3} each edge connecting $s\in L^{-1}_i (0)$ and $t\in L^{-1}_i(1)$ is parallel to $\vec e_i$,

\item\label{lemma:propertiesofL4} any fiber $L^{-1}(b)$, for $b\in \Bool{d}$, is contained in a facet of $P$.
\end{enumerate}
\end{lemma}

\skpt
\begin{proof}
\begin{enumerate}
\item When we switch a pair of successive left and right-leaning edges to a pair of successive right and left-leaning edges, either it does not change the orientation of the leaves or it just changes one left-leaning leaf into a right-leaning leaf.

\item This is straightforward from the definition of $\vec v$.

\item It is enough to prove it on one polytope $K_\omega$.
In this case, $t$ is a planar binary tree obtained from a planar binary tree $s$ by switching the $(i+1)^\text{th}$ leaf from left-leaning to right-leaning, which implies
\[
\overrightarrow{M(s, \omega)M(t, \omega)}=\omega_i\omega_{i+2}\vec e_i.
\]
\item Reading the sequence $(1, b_1, \ldots, b_{n_1-2}, 0)$ from left to right, there is at least one occurrence of $(1,0)$, say at place $i$ and $i+1$. Every face labeled by a forest of trees~$t$ satisfying $L(t)=b$ lies in the facet labeled by $\big(c_{n_1-1}\circ_i c_2, c_{n_2}, \ldots, c_{n_k}\big)$.\qedhere
\end{enumerate}
\end{proof}

\begin{lemma}\label{lem:FACE}
Let $F$ be a face of $P$ which contains a vertex $s$ such that $L_i(s) = 1$. For any $x\in \mathring F$, there exists $\varepsilon>0$ such that $x- \varepsilon \vec e_i\in P$.
\end{lemma}

\begin{proof}
It is enough to perform the proof for one Loday realization $P=\K_\omega$ that we endow with the linear form
$\psi(x)\coloneqq \ang{x, \vec e_i}$.
We claim that, for the projection $\id : P \to P$, the associated subcomplexes of lower and upper faces, introduced in \cref{subsec:PolySubDiv}, are given by
\[\Ve\big(\La^{\downarrow}(\tilde P)\big)=\{
M(t,\omega)\mid L_i(t)=0
\}
\quad
\text{and}
\quad
\Ve\big(\La^{\uparrow}(\tilde P)\big)=\{
M(t,\omega)\mid L_i(t)=1
\}.
\]
The intersection $L_i^{-1}(0)\cap \La^{\downarrow}(\tilde P)$
is nonempty, since it contains the vertex $\bm P$. Suppose that the polytopal complex
$\La^{\downarrow}(\tilde P)$ contains a vertex living in $L_i^{-1}(1)$. By~the connectivity of $\La^{\downarrow}(\tilde P)$, it admits an edge with vertices $s\in L_i^{-1}(0)$ and $t\in L_i^{-1}(1)$.
Point~\eqref{lemma:propertiesofL3} of \cref{lemma:propertiesofL} says that such an edge is parallel to $\vec e_i$, which contradicts the minimality of $\La^{\downarrow}(\tilde P)$.
This proves the inclusion
$\Ve\big(\La^{\downarrow}(\tilde P)\big)
\subset
\left\{
M(t,\omega)\mid L_i(t)=0
\right\}$.
The opposite inclusion is proved by the same argument as for
$\La^{\uparrow}(\tilde P)$
together with the fact that
the union $\La^{\downarrow}(\tilde P)\cup \La^{\uparrow}(\tilde P)$ contains all the vertices of $P$.
If not, a vertex~$x$, which is not contained in it, is neither maximal nor minimal with respect to $\vec e_i$, which contradicts the fact that $x$ is an extremal point.
This characterization of the polytopal complex of lower faces shows that any face $F$ of $P$ containing a vertex $s$ such that $L_i(s) = 1$ satisfies
$ F\not \subset\vert\La^{\downarrow}(\tilde P)\vert$, and thus
$\mathring F \cap \big\vert\La^{\downarrow}(\tilde P)\big\vert=\emptyset$, which concludes the proof.
\end{proof}

\begin{proposition}\label{prop:FirstExclusion}
Let $F$ and $G$ be two faces of $P$ of matching dimensions, \ie $\dim F+\dim G=d$. We consider
$s\coloneqq \tp F$ and $t\coloneqq \bm G$.
When $L(s)\not\leq L(t)$, we have $\mathring{F}\times \mathring{G}\cap \im\triangle = \emptyset$.
\end{proposition}

\begin{proof}
When $L(s)\not\leq L(t)$, there exists $i$ such that $L_i(s) = 1$ and $L_i(t) = 0$.
By~\cref{lem:FACE} and a dual version of it proved with the same arguments,
for every $x\in \mathring{F}$ and $y\in \mathring{G}$, there exists $\varepsilon>0$ such that $(x-\varepsilon \vec e_i, y+\varepsilon \vec e_i) \in P\times P$.
Suppose now that there exists $(x,y)\in \mathring{F}\times \mathring{G} \cap \im \triangle$.
Since $(x+y)/2 = ((x-\varepsilon\vec e_i)+(y+\varepsilon\vec e_i))/2$, the two points $(x, y)$ and $(x-\varepsilon\vec e_i, y+\varepsilon\vec e_i)$ lie in the same fiber of $\beta$. As we saw in the proof of
\cref{prop:DiaginPoly}, the point $(x,y)$ minimizes $\ang{x, \vec v}$ in the fiber of $\beta$.
However, the computation
\[
\ang{\vec v, x-\varepsilon \vec e_i}=\ang{\vec v, x}-\varepsilon\underbrace{\ang{\vec v, \vec e_i}}_{>0}
<\ang{\vec v, x}
\]
violates the definition of $\triangle$.
\end{proof}

\cref{prop:FirstExclusion} excludes
faces $F\times G$ of matching dimensions with $L(s)\not\leq L(t)$
from the image of $\triangle$.
Notice that, by Point~\eqref{lemma:propertiesofL1} of \cref{lemma:propertiesofL}, $L(s)\not\leq L(t)$ implies $s\not \leq t$.
In order to exclude the remaining case when $s\not\leq t$ and $L(s)\leq L(t)$, we prepare the following two lemmas.

\begin{lemma}\label{lem:UniqueFG}
Let $t\in \Tam{n_1}\times \cdots \times \Tam{n_k}$ be a forest of planar binary trees with a total of $r_t+k$ right-leaning leaves and $l_t+k$ left-leaning leaves. There exists a unique maximal (with respect to inclusion) face $F_t$ (\resp $G_t$) of $P$ such that \hbox{$\tp F\!=\!t$} (\resp $\bm G = t$). The dimensions of these faces are given by $\dim F_t =l_t$ and $\dim G_t = r_t$.
\end{lemma}

\begin{proof}
The cell $F_t$ (\resp $G_t$) can be obtained by collapsing all the left-leaning (\resp right-leaning) internal edges of all the trees of the forest $t$. One can see that, for any forest of planar binary trees, the number of left-leaning (\resp right-leaning) internal edges is equal to the number of right-leaning (\resp left-leaning) leaves minus~$k$.
\end{proof}

\begin{lemma}\label{lemma:Matching}
Let $F,G\subset P$ be a pair of faces of matching dimensions.
When $s\coloneqq \tp F$ and $t\coloneqq \bm G$ satisfy $L(s)\leq L(t)$, then $(F, G) = (F_s, G_t)$ and $L(s)=L(t)$.
\end{lemma}

\begin{proof}
By definition, $F\subset F_s$ and $G\subset G_t$, and thus $\dim F + \dim G\leq l_s + r_t\leq d$ by \cref{lem:UniqueFG}.
The top dimension assumption $\dim F+\dim G=d$ allows us to conclude that
$\dim F=l_s$, $\dim G=r_t$, $F=F_s$, $G=G_t$, and $L(s)=L(t)$.
\end{proof}

We can now conclude Step 1: we prove by induction on the dimension $d$ of $P$ that any pair $F, G\subset P$ of faces of matching dimensions with $s\not\leq t$ and $L(s)\leq L(t)$ satisfies
${F}\times {G}\not\subset \im\triangle$.
This is straightforward to check in dimensions $d=0$ and $d=1$.
In dimension $d$, let us suppose, to the contrary, that there exists such a pair $F,G$ of faces satisfying ${F}\times {G}\subset \im\triangle_{(P, \vec v)}$.
\cref{lemma:Matching} implies $L(s)=L(t)$.
In this case, both points~$s$ and~$t$ lie in a common facet $Q$ of $P$ by Point~\eqref{lemma:propertiesofL4} of~\cref{lemma:propertiesofL}.
By~\cref{Lemma:DiagOnFace}, the induced diagonal $\triangle_{(Q,\vec v)}$ on $Q$ is the restriction of ${\triangle_{(P, \vec v)}}$.
Let us consider $F'\coloneqq F\cap Q$ and $G'\coloneqq G\cap Q$. If $\dim F'+\dim G'> \dim Q$, then
$F'\times G'\not \subset \triangle_{(Q,\vec v)}$ by \cref{cor:DIMENSION}. When $\dim F'+\dim G'\leq \dim Q$, we consider any pair of faces $F'\subset F''$ and $G'\subset G''$ of $Q$ of matching dimensions: $\dim F''+\dim G''=\dim Q$.
They satisfy $r\coloneqq\tp F''\geq \tp F'=s$ and $u \coloneqq\bm G''\leq \bm G'=t$, which implies $r \not \leq u$.

We claim that ${F''}\times {G''}\not\subset \im\triangle_{(Q,\vec v)}$.
Since $Q$ is a facet of $P$, it is of type
$Q'\coloneqq K_{\bar \omega_1}\times K_{\tilde \omega_1}\times K_{\omega_2} \times \cdots \times K_{\omega_k}\cong Q$,
under the isomorphism $\Theta : \RR^N \cong \RR^N$ given by the block permutation of coordinates described in the proof of Point~\eqref{prop:PropertiesKLoday5} of \cref{prop:PropertiesKLoday}. The image of the orientation vector $\vec v=(v_1, \ldots, v_{n_1-1}, v_{n_1}, \ldots)$
under the inverse permutation of coordinates is equal to
\[\vec v':=\Theta^{-1}(\vec v)=(v_1, \ldots, v_{i-1}, v_{i+m-1}, \ldots, v_{n_1-1}, v_i, \ldots, v_{i+m-2}, v_{n_1}, \ldots, ),\]
so it well-orients $Q'$.
Therefore the isomorphism $\Theta$ intertwines the two diagonals $\triangle_{(Q,\vec v)}$ and $\triangle_{(Q',\vec v')}$. Let us denote by $r=(r_1, \ldots, r_k)$ and $u=(u_1, \ldots, u_k)$ the various planar trees composing the two forests. Under the notation of the proof of
Point~\eqref{prop:PropertiesKLoday5} of \cref{prop:PropertiesKLoday},
one can write the two planar trees
$r_1=\bar r_1\circ_i \tilde r_1$ and $u_1=\bar u_1\circ_i \tilde u_1$. This gives
\begin{align*}
r'&\coloneqq \tp \Theta^{-1}(F'')=(\bar r_1, \tilde r_1, r_2, \ldots, r_k)\\
\tag*{and} u' &\coloneqq \bm \Theta^{-1}(G'')=(\bar u_1, \tilde u_1, u_2, \ldots, u_k).
\end{align*}
We have $r'\not\leq u'$. Indeed, the condition $r\not \leq u$ implies that there exists $1\leq j\leq k$ such that $r_j \not \leq u_j$. If $j\neq 1$, then automatically $r'\not\leq u'$. If $j=1$, then either $\bar r_1 \not \leq \bar u_1$ or
$\tilde r_1 \not \leq \tilde u_1$, since $\circ_i$ preserves the Tamari order, and again $r'\not\leq u'$.
If $L(r')\not\leq L(u')$,
we conclude with \cref{prop:FirstExclusion}, otherwise we conclude our claim with the induction hypothesis.
Finally, notice that any cell appearing in $\im \triangle$ is contained in a product of cells of matching dimensions. The above argument shows that $F'\times G'$ cannot appear in $\im \triangle_{(Q,\vec v)}$. This concludes the first step.

\subsection{Second step: $\im\triangle_n\supset \bigcup F\times G$}
In this section, we prove that every matching pair $(F, G)$ of $\K_n$ satisfies $F\times G\subset \im\triangle_n$.
By Point \eqref{lemma:propertiesofL1} of \cref{lemma:propertiesofL} and by \cref{lemma:Matching} such pairs are of type
$(F_s, G_t)$ with $s\leq t$ and $L(s)=L(t)$.

\begin{proposition}\label{prop:FtGt}
For every $t\in \Tam{n}$, we have $F_t\times G_t\subset \im \triangle_n$.
\end{proposition}

\begin{proof}
By the pointwise formula $(t, t)\in \im\triangle_n$. But $F_t\times G_t$ is the only cell coming from a matching pair that can contain $(t, t)$ by \cref{Subsec:Subset}.
\end{proof}

\begin{lemma}\label{lem:FacetGuGt}
Let $t\prec u$ be an edge of $\K_n$ satisfying $L(t) = L(u)$.
\begin{enumerate}
\item\label{lem:FacetGuGt1} The cell $G_t\cap G_u$ is a facet of $G_t$ and $G_u$.

\item\label{lem:FacetGuGt2} The cell $G_t\cap G_u$ is not of type $G_v$, for $v\in \Tam{n}$.

\item\label{lem:FacetGuGt3} If $G_t\cap G_u$ is a facet of $G_v$, then $v=t$ or $v=u$.
\end{enumerate}
\end{lemma}

\skpt
\begin{proof}
\begin{enumerate}
\item Let the planar binary tree $u$ be obtained from $t$ by switching
a pair $(f,e)$ of successive left and right-leaning edges
to a pair $(f', e')$ of successive
right and left-leaning edges:
\[
t=\ {\vcenter{\hbox{
\begin{tikzpicture}[yscale=0.5,xscale=0.5]
\draw[thick] (0,-1)--(0,0) -- (-2,2);
\draw[thick] (-1,1)--(0,2) ;
\draw[thick] (0,0)--(1,1) ;
\draw [fill] (0,-1) circle [radius=0.014];
\draw [fill] (-2,2) circle [radius=0.014];
\draw [fill] (1,1) circle [radius=0.014];
\draw [fill] (0,2) circle [radius=0.014];

\draw (-2,2) node[above] {$t_1$};
\draw (0,2) node[above] {$t_2$};
\draw (1,1) node[above] {$t_3$};
\draw (0,-1) node[below] {$t_4$};
\draw (-0.7,0.35) node[above right] {\scalebox{0.8}{$e$}};
\draw (-0.25,1.25) node[above left] {\scalebox{0.8}{$f$}};
\end{tikzpicture}}}}
\qquad \text{and} \qquad
u=\ {\vcenter{\hbox{
\begin{tikzpicture}[yscale=0.5,xscale=0.5]
\draw[thick] (0,-1)--(0,0) -- (2,2);
\draw[thick] (1,1)--(0,2) ;
\draw[thick] (0,0)--(-1,1) ;
\draw [fill] (0,-1) circle [radius=0.014];
\draw [fill] (2,2) circle [radius=0.014];
\draw [fill] (-1,1) circle [radius=0.014];
\draw [fill] (0,2) circle [radius=0.014];

\draw (2,2) node[above] {$t_3$};
\draw (0,2) node[above] {$t_2$};
\draw (-1,1) node[above] {$t_1$};
\draw (0,-1) node[below] {$t_4$};
\draw (0.8,0.3) node[above left] {\scalebox{0.8}{$e'$}};
\draw (0.4,1.35) node[above right] {\scalebox{0.8}{$f'$}};
\end{tikzpicture}}}}.
\]
By \cref{lem:UniqueFG}, the cell $G_t\cap G_u$ corresponds to a tree $s$ obtained by collapsing all the right-leaning internal edges of $u$ except $f'$.
The cell $G_u$ corresponds to the tree
$s / f'$ obtained from $s$ by contracting the edge $f'$.
The cell $G_t$ corresponds to the tree $s/ e'$.

\item Notice first that a cell $G$ is of type $G_v$, for $v\in \Tam{n}$ if and only if it is labeled by a planar tree
having only left-leaning internal edges.
This is not the case of the tree $s$, since it carries the right-leaning internal edge $f'$.

\item In this case, the labeling tree $r$ of $G_v$ is obtained from $s$ by contracting one internal edge. Since $r$ should not contain any right-leaning internal edge, it can only be obtained in two ways: by contracting $e'$ or $f'$. In the former case $v=t$ and in the latter case $v=u$.\qedhere
\end{enumerate}
\end{proof}

\begin{proposition}
If $s\leq t$ and $L(s) = L(t)$ then $F_s\times G_t\subset \im\triangle_n$.
\end{proposition}

\begin{proof}
We fix $s$ and we consider the sub-lattice $\Tam{n}^{\geq s}$ of elements greater than~$s$.
Let us prove by induction on the elements $t$ of $\Tam{n}^{\geq s}\cap L^{-1}(L(s))$ that
$F_s\times G_t\subset \im\triangle_n$.
The base case is $s=t$ and this is done in the above \cref{prop:FtGt}.
Suppose the conclusion holds for $t$ and let $t\prec u$ be an edge satisfying $L(t)= L(u)$.
Then $H \coloneqq F_s\times G_t\cap F_s\times G_u = F_s\times (G_t\cap G_u)$ is a facet of $F_s\times G_t$
by Point~\eqref{lem:FacetGuGt1} of \cref{lem:FacetGuGt} and it does not project into the boundary of $\K_n$ via $\beta$.
Since $F_s\times G_t$ lies in $\im\triangle_n$, it is attached along $H$ to a cell
$F_r\times G_v$ by \cref{Subsec:Subset}.
We claim that $r=s$, $v=u$.
Since $H=F_s\times (G_t\cap G_u)$ is a facet of $F_r\times G_v$, there are two options: either $F_s=F_r$ or $G_t\cap G_u=G_v$. The latter case is contradicted by Point~\eqref{lem:FacetGuGt2} of \cref{lem:FacetGuGt}.
Therefore, $G_t\cap G_u$ is a facet of $G_v$, and we conclude by
Point~\eqref{lem:FacetGuGt3} of \cref{lem:FacetGuGt}.
\end{proof}

This concludes the proof of Step 2.

\begin{thebibliography}{99}
\bibitem{AguiarArdila17} M. Aguiar \& F. Ardila - “Hopf monoids and generalized permutahedra” , 2017
\bibitem{Amorim17} L. Amorim - “The Künneth theorem for the Fukaya algebra of a product of Lagrangians” , Internat. J. Math. 28 (2017) no. 4, article ID 1750026, 38 pages
\bibitem{Brown59} E. H. Brown - “Twisted tensor products. I” , Ann. of Math. (2) 69 (1959), p. 223-246
\bibitem{BilleraSturmfels92} L. J. Billera \& B. Sturmfels - “Fiber polytopes” , Ann. of Math. (2) 135 (1992) no. 3, p. 527-549
\bibitem{BoardmanVogt73} J. M. Boardman \& R. M. Vogt - Homotopy invariant algebraic structures on topological spaces , Lect. Notes in Math. , vol. 347 , Springer-Verlag, Berlin, 1973
\bibitem{CFZ02} F. Chapoton, S. Fomin \& A. Zelevinsky - “Polytopal realizations of generalized associahedra” , Canad. J. Math. 45 (2002) no. 4, p. 537-566
\bibitem{CSZ15} C. Ceballos, F. Santos \& G. M. Ziegler - “Many non-equivalent realizations of the associahedron” , Combinatorica 35 (2015) no. 5, p. 513-551
\bibitem{EilenbergMacLane54} S. Eilenberg \& S. Mac Lane - “On the groups $H(\Pi ,n)$. II. Methods of computation” , Ann. of Math. (2) 60 (1954), p. 49-139
\bibitem{EilenbergZilber53} S. Eilenberg \& J. A. Zilber - “On products of complexes” , Amer. J. Math. 75 (1953), p. 200-204
\bibitem{Forcey08} S. Forcey - “Convex hull realizations of the multiplihedra” , Topology Appl. 156 (2008) no. 2, p. 326-347
\bibitem{GKZ08} I. M. Gelfand, M. M. Kapranov \& A. V. Zelevinsky - Discriminants, resultants and multidimensional determinants , Modern Birkhäuser Classics , Birkhäuser Boston, Inc., Boston, MA, 2008
\bibitem{GaberdielZwiebach97} M. R. Gaberdiel \& B. Zwiebach - “Tensor constructions of open string theories. I. Foundations” , Nuclear Phys. B 505 (1997) no. 3, p. 569-624
\bibitem{Loday04a} J.-L. Loday - “Realization of the Stasheff polytope” , Arch. Math. (Basel) 83 (2004) no. 3, p. 267-278
\bibitem{LodayVallette12} J.-L. Loday \& B. Vallette - Algebraic operads , Grundlehren Math. Wiss. , vol. 346 , Springer-Verlag, Berlin, 2012
\bibitem{May72} J. May - The geometry of iterated loop spaces , Lect. Notes in Math. , vol. 271 , Springer-Verlag, Berlin, 1972
\bibitem{MarklShnider06} M. Markl \& S. Shnider - “Associahedra, cellular $W$-construction and products of $A_\infty $-algebras” , Trans. Amer. Math. Soc. 358 (2006) no. 6, p. 2353-2372
\bibitem{MSS} M. Markl, S. Shnider \& J. D. Stasheff - Operads in algebra, topology and physics , Math. Surveys and Monographs , vol. 96 , American Mathematical Society, Providence, RI, 2002
\bibitem{Proute11} A. Prouté - “$A_\infty $-structures. Modèles minimaux de Baues-Lemaire et Kadeishvili et homologie des fibrations” , Repr. Theory Appl. Categ. (2011) no. 21, p. 1-99, Reprint of the 1986 original, With a preface to the reprint by Jean-Louis Loday
\bibitem{Seidel08} P. Seidel - Fukaya categories and Picard-Lefschetz theory , Zürich Lectures in Advanced Math. , European Mathematical Society, Zürich, 2008
\bibitem{Serre51} J.-P. Serre - “Homologie singulière des espaces fibrés. Applications” , Ann. of Math. (2) 54 (1951), p. 425-505
\bibitem{StasheffShnider97} J. D. Stasheff \& S. Shnider - “From operads to “physically” inspired theories (Appendix B)” , in Operads: Proceedings of Renaissance Conferences (Hartford, CT/Luminy, 1995) , Contemp. Math. , vol. 202 , American Mathematical Society, Providence, RI, 1997, p. 53-81
\bibitem{Stasheff63} J. D. Stasheff - “Homotopy associativity of $H$-spaces. I \& II” , Trans. Amer. Math. Soc. 108 (1963), p. 275-292 \& 293–312
\bibitem{Stasheff70} J. D. Stasheff - $H$-spaces from a homotopy point of view , Lect. Notes in Math. , vol. 161 , Springer-Verlag, Berlin, 1970
\bibitem{Steenrod47} N. E. Steenrod - “Products of cocycles and extensions of mappings” , Ann. of Math. (2) 48 (1947), p. 290-320
\bibitem{SaneblidzeUmble04} S. Saneblidze \& R. Umble - “Diagonals on the permutahedra, multiplihedra and associahedra” , Homology Homotopy Appl. 6 (2004) no. 1, p. 363-411
\bibitem{Tamari51} D. Tamari - Monoïdes préordonnés et chaînes de Malcev , Thèse de Mathématique, Paris, 1951
\bibitem{Ziegler95} G. M. Ziegler - Lectures on polytopes , Graduate Texts in Math. , vol. 152 , Springer-Verlag, New York, 1995
\end{thebibliography}
\end{document}
